\subsubsection{Quantum Oracle \texorpdfstring{$U_f$}{Uf}}\label{ex:quantum-oracle-uf}

In the classical version of the problem, the oracle receives one bit $x$ and returns $f(x)$, where $f$ is a function that maps $\{0,1\}$ to $\{0,1\}$. Since Deutsch's algorithm is a \textbf{quantum algorithm}, we need a \hl{quantum circuit that performs the same computation.} A first idea would be to define a quantum operator satisfying:
\begin{equation*}
    \ket{x} \mapsto \ket{f(x)}
\end{equation*}
It means that the oracle would take a qubit in state $\ket{x}$ and return a qubit in state $\ket{f(x)}$. Unfortunately, this construction is \textbf{not valid in general}.

\highspace
\textcolor{Green3}{\faIcon{question-circle} \textbf{Why doesn't this work?}} Every operation performed by a quantum computer must be represented by a \textbf{unitary operator}. One fundamental property of unitary operators is that they are \textbf{reversible}. In other words, every output state must uniquely determine the corresponding input state. Mathematically:
\begin{equation*}
    U^{-1}U = UU^{-1} = I
\end{equation*}
Therefore, no two different input states may evolve into the same output state.

\highspace
If we're not sure why this doesn't work, consider the example of a constant function $f(x) = 0$. Using the naive mapping:
\begin{equation*}
    \ket{x} \mapsto \ket{f(x)}
\end{equation*}
\textcolor{Red2}{\faIcon{exclamation-triangle}} We obtain:
\begin{align*}
    \ket{0} &\mapsto \ket{f(0)} = \ket{0} \\
    \ket{1} &\mapsto \ket{f(1)} = \ket{0}
\end{align*}
But this is \textbf{not reversible}, since both $\ket{0}$ and $\ket{1}$ are mapped to the same output state $\ket{0}$. Therefore, this \textbf{mapping is not a valid quantum operation}.

\highspace
\begin{flushleft}
    \textcolor{Green3}{\faIcon{check-circle} \textbf{Preserving the input}}
\end{flushleft}
The solution is remarkably simple. Instead of replacing the input by the function value, we \textbf{keep the original input unchanged} (\emph{data qubit}) and \textbf{store the result in a second qubit} (\emph{computation qubit}). The oracle therefore acts on \textbf{two qubits} instead of one. The first qubit stores the function input $x$, while the second qubit stores an auxiliary value, $y$. The oracle computes the function by modifying only the second qubit (the auxiliary qubit).

\highspace
The standard \definitionWithSpecificIndex{Quantum Oracle}{Quantum Oracle}{Oracle} $U_f$ is defined as follows:
\begin{equation}
    U_f\ket{x, y} = \ket{x, y \oplus f(x)}
\end{equation}
Where:
\begin{itemize}
    \item $x$ is the input bit to the function $f$. It is called \definition{Data Qubit} and \textbf{contains the function input} $\ket{x}$ and \textbf{remains unchanged} throught the oracle evaluation.

    \item $y$ is an auxiliary (or computation) bit that is used to \textbf{store the result of the function} evaluation. It is called \definition{Computation Qubit} and stores the auxiliary value $\ket{y}$, which becomes $\ket{y \oplus f(x)}$ after the oracle evaluation. It is called also \definition{Ancilla Qubit} or \definition{Target Qubit}.

    \item $\oplus$ denotes addition modulo 2 (XOR operation).\footnote{%
        The XOR operation is defined as follows:
        \begin{equation*}
            0 \oplus 0 = 0, \quad 0 \oplus 1 = 1, \quad 1 \oplus 0 = 1, \quad 1 \oplus 1 = 0
        \end{equation*}%
    }
\end{itemize}
\textcolor{Green3}{\faIcon{question-circle} \textbf{Why does XOR make the oracle reversible?}} The XOR operation has a crucial property:
\begin{equation*}
    \left(a \oplus b\right) \oplus b = a
\end{equation*}
Applying the same XOR twice restores the original value. Indeed:
\begin{align*}
    U_f^2 \ket{x,y} &= U_f \ket{x, y \oplus f(x)} \\
    &= \ket{x, (y \oplus f(x)) \oplus f(x)} \\
    &= \ket{x, y}
\end{align*}
Therefore:
\begin{equation*}
    U_f^{-1} = U_f
\end{equation*}
Meaning that the \textbf{oracle is its own inverse}. The transformation is reversible for \textbf{every} Boolean function, including the constant ones.

\highspace
\begin{flushleft}
    \textcolor{Green3}{\faIcon{tools} \textbf{Oracle implementations for the four functions}}
    \phantomsection\label{ex:oracle-implementations-for-the-four-functions}
\end{flushleft}
The four possible Boolean functions correspond to four different quantum circuits:
\begin{itemize}
    \item \hl{$f(x)=0$}, the oracle action is:
    \begin{equation*}
        \ket{x,y} \mapsto \ket{x,y}
    \end{equation*}
    The circuit implementation is an \textbf{identity gate}:
    \begin{center}
        \begin{quantikz}
            \lstick{$\ket{x}$} & & \\
            \lstick{$\ket{y}$} & & 
        \end{quantikz}
    \end{center}

    \item \hl{$f(x)=1$}, the oracle action is:
    \begin{equation*}
        \ket{x,y} \mapsto \ket{x,y \oplus 1}
    \end{equation*}
    The circuit implementation is a \textbf{Pauli-$X$ gate on the computation qubit}:
    \begin{center}
        \begin{quantikz}
            \lstick{$\ket{x}$} & & \\
            \lstick{$\ket{y}$} & \gate{X} & 
        \end{quantikz}
    \end{center}

    \item \hl{$f(x)=x$}, the oracle action is:
    \begin{equation*}
        \ket{x,y} \mapsto \ket{x,y \oplus x}
    \end{equation*}
    The circuit implementation is a \textbf{Controlled-NOT gate}, with the \textbf{data qubit as control} and the \textbf{computation qubit as target}:
    \begin{center}
        \begin{quantikz}
            \lstick{$\ket{x}$} & \ctrl{1} & \\
            \lstick{$\ket{y}$} & \targ{} & 
        \end{quantikz}
    \end{center}

    \item \hl{$f(x)=\overline{x}$}, the oracle action is:
    \begin{equation*}
        \ket{x,y} \mapsto \ket{x,y \oplus \overline{x}}
    \end{equation*}
    Since:
    \begin{equation*}
        \overline{x} = 1 \oplus x
    \end{equation*}
    We can write:
    \begin{equation*}
        y \oplus \overline{x} = y \oplus 1 \oplus x
    \end{equation*}
    The circuit implementation is:
    \begin{enumerate}
        \item Apply a CNOT from $x$ to $y$;
        \item Apply a Pauli-$X$ gate to the \textbf{computation qubit} $y$.
    \end{enumerate}
    The order may also reversed because an $X$ on the target commutes with this CNOT:
    \begin{equation*}
        CX\left(I \otimes X\right) = \left(I \otimes X\right) CX
    \end{equation*}
    \begin{center}
        \begin{quantikz}
            \lstick{$\ket{x}$} & \ctrl{1} &          & \\
            \lstick{$\ket{y}$} & \targ{}  & \gate{X} & 
        \end{quantikz}
        =
        \begin{quantikz}
            \lstick{$\ket{x}$} &          & \ctrl{1} & \\
            \lstick{$\ket{y}$} & \gate{X} & \targ{}  &
        \end{quantikz}
    \end{center}
\end{itemize}

\noindent
Notice an interesting pattern:
\begin{itemize}
    \item The \textbf{constant} functions require \textbf{no controlled operation};
    \item The \textbf{balanced} functions always involve a \textbf{controlled-NOT gate}.
\end{itemize}
At this point this may seem like a coincidence. In the next sections we will see that Deutsch's algorithm exploits this difference using \textbf{phase kickback}, allowing the presence or absence of the controlled operation to be detected with a single oracle query.